\section{Identification and upper-bound proofs}

The upper bound in \cref{obj:thm:adaptive-minimax} combines causal identification with control of approximation error, conditional stochastic error, and exceptional samples. The auxiliary statements below follow that order: identification fixes the population target, cube smoothness admits a controlled global extension, the global overlap tail supplies an ordered information profile, and the template and mesh bounds translate that profile into stable response estimation.

\subsection{Identification of the continuous response}

Under the causal conditions in \cref{obj:def:model}, the treated conditional mean identifies the causal response almost everywhere. The density and smoothness restrictions then determine its continuous representative throughout the covariate cube. The following statement records both levels of identification.

\begin{lemmav}{lem:causal-identification}[Identification of the treated response]
Let \(P\in\mathcal P_\gamma\) satisfy the parameter restrictions and causal model conditions of \cref{obj:def:model}, with \(\gamma>1\). Write \(P_X\) for its covariate marginal and \(\mathcal X=[0,1]^d\) for the covariate cube. Then
\[
P_X\{x:e(x)=0\}=0,
\qquad
\mu_1(x)=\mathbb E_P[Y\mid A=1,X=x]
\quad\text{for }P_X\text{-almost every }x.
\]
Moreover, every function \(g:\mathbb R^d\to\mathbb R\) that is continuous on \(\mathcal X\) and satisfies
\[
g(x)=\mathbb E_P[Y\mid A=1,X=x]
\quad\text{for }P_X\text{-almost every }x
\]
obeys
\[
g(x)=\mu_1(x)
\qquad\text{for every }x\in\mathcal X.
\]
\end{lemmav}

The pointwise and spatial-supremum losses in \cref{obj:def:risks} therefore concern a uniquely identified continuous response, including at boundary points. This separates the identification of the target from the statistical difficulty created by small treated probabilities.

\subsection{Global extension of cube smoothness}

The intrinsic smoothness convention in \cref{obj:def:holder} specifies coordinate derivatives through their continuous boundary traces. For an integer derivative order \(m\ge0\) and a top-derivative Hölder exponent \(s\in(0,1]\), the following statement supplies a global representative with controlled Fréchet derivatives. In the response class, these orders are \(m=\lceil\beta\rceil-1\) and \(s=\beta-m\).

\begin{lemmav}{lem:global-holder-extension}[Global Hölder extension bounds]
Fix an integer \(d\ge0\), an integer \(m\ge0\), and an exponent \(0<s\le1\). There exists a constant \(K_{\mathrm{ext}}>0\), depending only on \(d\), \(m\), and \(s\), with the following property. Let \(L\ge0\) and let \(g:[-1,1]^d\to\mathbb R\) satisfy:
\begin{itemize}
\item \textbf{(Intrinsic smoothness.)} The function \(g\) belongs to \(C^m([-1,1]^d)\) in the intrinsic closed-cube sense of \cref{obj:def:holder}, transported to \([-1,1]^d\). Its coordinate derivatives \(D^\alpha g\), for \(|\alpha|\le m\), exist on the interior and admit continuous traces on the closed cube.
\item \textbf{(Derivative bounds.)} With derivatives evaluated through their continuous traces,
\[
\sup_{z\in[-1,1]^d}|D^\alpha g(z)|\le L
\qquad\text{for every }|\alpha|\le m.
\]
\item \textbf{(Hölder continuity.)} For every \(|\alpha|=m\) and all \(z,z'\in[-1,1]^d\),
\[
|D^\alpha g(z)-D^\alpha g(z')|
\le L\|z-z'\|_\infty^{s}.
\]
\end{itemize}
Then there exists \(g_{\mathrm{ext}}\in C^m(\mathbb R^d)\) such that
\[
g_{\mathrm{ext}}(z)=g(z)
\qquad\text{for every }z\in[-1,1]^d,
\]
\[
\sup_{z\in\mathbb R^d}
\|\mathrm D^{k}g_{\mathrm{ext}}(z)\|_{\mathrm{op}}
\le K_{\mathrm{ext}}L
\qquad\text{for every }0\le k\le m,
\]
and
\[
\|\mathrm D^{m}g_{\mathrm{ext}}(z)
-\mathrm D^{m}g_{\mathrm{ext}}(z')\|_{\mathrm{op}}
\le K_{\mathrm{ext}}L\|z-z'\|_\infty^{s}
\qquad\text{for all }z,z'\in\mathbb R^d.
\]
Here \(\mathrm D^{k}\) denotes the Fréchet derivative of order \(k\), with order zero interpreted as the function itself, and \(\|\cdot\|_{\mathrm{op}}\) denotes the operator norm of a \(k\)-linear form on \(\mathbb R^d\) equipped with the supremum norm \(\|\cdot\|_\infty\).
\end{lemmav}

The extension preserves every value on the closed cube and converts the coordinate-derivative restrictions into global operator-norm bounds. Its smoothness factor \(K_{\mathrm{ext}}\) depends on dimension and smoothness order, uniformly over functions satisfying the displayed radius bound. This supplies a controlled global representation of intrinsic cube smoothness while preserving the response values relevant to approximation.

\subsection{The ordered treated-mass profile}

For a dyadic mesh \(h\), let \(q_{Q\ell}\) denote the population treated probability of template microcell \(E_{Q\ell}\) in partition cube \(Q\). The quantity \(q_Q\) is the smallest of these microcell probabilities within \(Q\), and \(q_{(k)}\) is the \(k\)th smallest entry among the cube-level quantities \(q_Q\). Ties may be ordered arbitrarily. These masses concern the scaled microcells of \cref{obj:def:template}, rather than the total treated probability of each partition cube.

The global-tail restriction in \cref{obj:ass:global-tail} controls treated probability on every measurable covariate region. Its implication for the ordered weakest-cell masses supplies the information profile used in the stochastic analysis.

\begin{lemmav}{lem:ordered-mass}[Uniform ordered treated masses]
Fix the model parameters and an overlap interval satisfying
\begin{itemize}
\item \textbf{(Parameter domain.)} \(d\in\mathbb N\), \(d\ge1\), \(\beta>1\), \(B>0\), \(L>0\), \(C\ge1\), and \(0<c_f\le1\).
\item \textbf{(Overlap range.)} \(1<\gamma_{\min}<\gamma_{\max}\), with \(\Gamma=[\gamma_{\min},\gamma_{\max}]\).
\end{itemize}
There exists \(\kappa_\Gamma>0\), depending only on these fixed parameters and interval endpoints, such that the following two bounds hold for every \(\gamma\in\Gamma\) and every causal law \(P\in\mathcal P_\gamma\) satisfying \cref{obj:def:model}.

For every measurable \(E\subseteq\mathcal X=[0,1]^d\), writing \(u=P(X\in E)\),
\[
P(A=1,X\in E)
\ge
\frac{\gamma-1}{\gamma}\,
C^{-1/(\gamma-1)}u^{\gamma/(\gamma-1)}.
\]

For each \(j\in\mathbb N\), set
\[
h=2^{-j},\qquad
m=\lceil\beta\rceil-1,\qquad
D_\gamma=\frac{d\gamma}{\gamma-1}.
\]
Let \(\mathcal Q_h\) be the dyadic cube partition of \(\mathcal X\), with \(|\mathcal Q_h|=(2^j)^d\). Using the degree-\(m\) polynomial template microcells \(E_{Q\ell}\) within each \(Q\in\mathcal Q_h\), define their treated masses and each cube's weakest mass by
\[
q_{Q\ell}=P(A=1,X\in E_{Q\ell}),
\qquad
q_Q=\min_{1\le\ell\le J}q_{Q\ell},
\]
where \(J\) is the number of template microcells. Let \(q_{(k)}\) be the \(k\)th entry of the list of these \(q_Q\)'s sorted in nondecreasing order. Then, for every integer \(1\le k\le(2^j)^d\),
\[
q_{(k)}
\ge
\kappa_\Gamma h^{D_\gamma}k^{1/(\gamma-1)}.
\]
\end{lemmav}

The rank factor quantifies how information improves across the ordered cubes. At a fixed mesh, the weakest treated mass grows at least polynomially with rank, so the corresponding variance envelope decreases with rank. The constant \(\kappa_\Gamma\) is common to the stated compact overlap range; this uniformity carries the mass comparison into the adaptive upper bound.

\subsection{Deterministic conditioning of the template}

Write \(U(x)\) for the vector of monomials of total degree at most \(m\), with \(m=\lceil\beta\rceil-1\) in the statistical application. The tensor nodes \(v_\ell\), reference microcubes \(V_\ell\), and reference Gram matrix \(G_0\) are those of \cref{obj:def:template}. Its least eigenvalue is denoted by \(\lambda_0\). The empirical matrix \(\widehat G_Q\) assigns equal weight to the microcells and averages within each cell over its treated observations. The next statement gives the conditioning bound for every feasible sampled configuration.

\begin{lemmav}{lem:template-conditioning}[Uniform equal-cell Gram conditioning]
Let \(d,m,n\) be nonnegative integers, let \(\beta\in\mathbb R\), and let \(\mathcal D_n=((X_i,A_i,Y_i))_{i=1}^n\) be any sample with \(X_i\in\mathbb R^d\), \(A_i\in\{0,1\}\), and \(Y_i\in\mathbb R\). Write \(U(x)=(x^\alpha)_{|\alpha|\le m}\) for the vector of total-degree-\(m\) monomials. For the polynomial template's tensor nodes \(v_\ell\) and reference microcubes \(V_\ell\), indexed by \(\ell\in\{0,\ldots,m\}^d\), define
\[
J=(m+1)^d,\qquad
G_0=\frac1J\sum_{\ell\in\{0,\ldots,m\}^d}
U(v_\ell)U(v_\ell)^\top,\qquad
\lambda_0=\inf_{\|b\|_2=1}b^\top G_0b.
\]
Here \(\|b\|_2=(\sum_{|\alpha|\le m}b_\alpha^2)^{1/2}\).

Let \(h=2^{-j}\), with \(j\) a nonnegative integer, be a feasible candidate dyadic mesh: \(h\in\mathcal H_n\), the candidate mesh collection, and
\[
N_{Q\ell}\ge h^{-2\beta}
\quad\text{for every }Q\in\mathcal Q_h
\text{ and }\ell\in\{0,\ldots,m\}^d,
\]
where \(\mathcal Q_h\) is the dyadic cube partition at mesh \(h\). For a cube \(Q\) with corner \(a_Q\), the microcell and its treated count are
\[
E_{Q\ell}=a_Q+hV_\ell,\qquad
N_{Q\ell}=\sum_{i=1}^n
\mathbf1\{A_i=1,\ X_i\in E_{Q\ell}\}.
\]
Define the equal-cell empirical Gram matrix by
\[
\widehat G_Q
=\frac1J\sum_{\ell\in\{0,\ldots,m\}^d}
\frac1{N_{Q\ell}}
\sum_{i=1}^n
\mathbf1\{A_i=1,\ X_i\in E_{Q\ell}\}
U\!\left(\frac{X_i-a_Q}{h}\right)
U\!\left(\frac{X_i-a_Q}{h}\right)^\top.
\]
Then \(\lambda_0>0\), and for every \(Q\in\mathcal Q_h\),
\[
\frac{\lambda_0}{2}\|b\|_2^2
\le b^\top\widehat G_Qb
\quad\text{for every }b=(b_\alpha)_{|\alpha|\le m},
\qquad
\det(\widehat G_Q)\ne0,
\qquad
\|\widehat G_Q^{-1}\|_{\mathrm{op}}
\le\frac2{\lambda_0},
\]
where \(\|M\|_{\mathrm{op}}=\sup_{\|b\|_2=1}\|Mb\|_2\).
\end{lemmav}

Equal weighting preserves the reference template's conditioning even when the treated counts differ across microcells. The bound is deterministic once feasibility holds, and it controls the amplification of both approximation residuals and outcome fluctuations in the fitted coefficients. Its geometric requirement is supplied by the occupied template microcells and their prescribed widths.

\subsection{Selected-mesh concentration and maximal error}

Let \(\mathcal D_n=((X_i,A_i,Y_i))_{i=1}^n\) be the observed sample. Conditional expectations below condition on the sampled design \((X_i,A_i)_{i=1}^n\), comprising the covariates and treatment indicators. Under this conditioning, the selected mesh of \cref{obj:def:mesh-selector} is fixed. The coefficient vector \(\widehat c_Q\) is the equal-cell polynomial fit used in \cref{obj:def:estimator}; centering it by its conditional expectation isolates the stochastic coefficient error.

The oracle mesh is \(h_{\gamma,n}=n^{-1/(2\beta+D_\gamma)}\), and the oracle response-recovery scale is \(r_{\gamma,n}=h_{\gamma,n}^{\beta}\). The sample event \(\mathcal E_n\) below controls both selection and the comparison of empirical counts with population treated masses. For the accompanying maximal inequality, \(Z_1,\ldots,Z_N\) denotes a generic family of centered real variables under a probability law \(\mu\). The positive quantities \(a\) and \(b\) specify a common scale cap and an ordered decay scale, respectively.

\begin{lemmav}{lem:selected-mesh}[Selected mesh and coefficient envelope]
Fix the dimension \(d\), smoothness parameter \(\beta\), model constants \(B,L,C,c_f\), and adaptation endpoints \(\gamma_{\min},\gamma_{\max}\), satisfying
\begin{itemize}
\item \textbf{(Dimension and smoothness.)} \(d\in\mathbb N\), \(d\ge1\), and \(\beta>1\).
\item \textbf{(Model constants.)} \(B>0\), \(L>0\), \(C\ge1\), and \(0<c_f\le1\).
\item \textbf{(Adaptation interval.)} \(1<\gamma_{\min}<\gamma_{\max}\), with
\[
\Gamma=[\gamma_{\min},\gamma_{\max}].
\]
\end{itemize}
There exist constants \(K,K'>0\) and \(n_0\in\mathbb N\) such that the following holds uniformly for every \(n\ge n_0\), \(\gamma\in\Gamma\), and observed law \(P\in\mathcal P_\gamma\), with the displayed fixed model constants. Let the sample have the iid law \(P^{\otimes n}\) specified in \cref{obj:ass:iid}, and use the mesh selector of \cref{obj:def:mesh-selector}. Write
\[
m=\lceil\beta\rceil-1,
\qquad
h_{\gamma,n}=n^{-1/(2\beta+D_\gamma)},
\]
where \(D_\gamma\) is the effective dimension governing the model's overlap-dependent rate, and let \(r_{\gamma,n}\) denote its oracle response-recovery risk scale.

For each partition cube \(Q\) and template microcell \(E_{Q\ell}\), define the population treated mass, its empirical count, and the weakest treated microcell mass by
\[
q_{Q\ell}=P\{A=1,\ X\in E_{Q\ell}\},
\qquad
N_{Q\ell}=\sum_{i=1}^n
\mathbf1\{A_i=1,\ X_i\in E_{Q\ell}\},
\qquad
q_Q=\min_\ell q_{Q\ell}.
\]
There is a measurable sample event \(\mathcal E_n\) such that
\[
P^{\otimes n}(\mathcal E_n^c)\le n^{-2}.
\]
For every sample in \(\mathcal E_n\), the selected mesh satisfies
\[
0<\widehat h\le K h_{\gamma,n},
\qquad
N_{Q\ell}\ge \frac{nq_{Q\ell}}5
\quad
\text{for every }Q\in\mathcal Q_{\widehat h}
\text{ and every }\ell.
\]

Order the selected cubes by increasing \(q_Q\), with ranks starting at \(1\). Conditional on the sampled design \((X_i,A_i)_{i=1}^n\), evaluate all coefficient estimators at the selected mesh of the given sample. For the cube of rank \(k\), every monomial coefficient indexed by a multiindex \(\alpha\) of total degree at most \(m\), and every \(t\in\mathbb R\),
\[
\begin{aligned}
&\mathbb E\!\left[
\exp\!\left\{
t\left(
\widehat c_{Q,\alpha}
-\mathbb E[\widehat c_{Q,\alpha}\mid(X_i,A_i)_{i=1}^n]
\right)
\right\}
\,\middle|\,(X_i,A_i)_{i=1}^n
\right]
\\
&\hspace{2em}\le
\exp\!\left\{
\frac{Kt^2}{2}
\min\!\left(
\widehat h^{2\beta},
n^{-1}\widehat h^{-D_\gamma}k^{-1/(\gamma-1)}
\right)
\right\}.
\end{aligned}
\]
These conditional exponential moments are integrable for every coefficient and every \(t\), and the maximum coefficient-vector error is conditionally integrable. With \(\|\cdot\|_\infty\) denoting the maximum absolute coordinate over those monomial indices, its conditional expectation satisfies
\[
\begin{aligned}
&\mathbb E\!\left[
\max_{Q\in\mathcal Q_{\widehat h}}
\left\|
\widehat c_Q-
\mathbb E[\widehat c_Q\mid(X_i,A_i)_{i=1}^n]
\right\|_\infty
\,\middle|\,(X_i,A_i)_{i=1}^n
\right]
\\
&\hspace{2em}\le
K\widehat h^\beta
\sqrt{1+\max\{0,\log(h_{\gamma,n}/\widehat h)\}}
\le K'r_{\gamma,n}.
\end{aligned}
\]

Moreover, for every exponent \(\alpha>0\), there exists \(K_\alpha>0\) such that the following maximal inequality holds. Let \(\mu\) be a probability law on a measurable space, let \(N\in\mathbb N\), and let \(Z_1,\ldots,Z_N\) be real random variables satisfying
\begin{itemize}
\item \textbf{(Scales.)} \(a>0\) and \(b>0\).
\item \textbf{(Exponential integrability.)} For every \(1\le k\le N\) and \(t\in\mathbb R\), \(\exp(tZ_k)\) is integrable under \(\mu\).
\item \textbf{(Moment envelope.)} For every \(1\le k\le N\) and \(t\in\mathbb R\),
\[
\mathbb E_\mu[\exp(tZ_k)]
\le
\exp\!\left\{
\frac{t^2}{2}\min(a^2,b^2k^{-\alpha})
\right\}.
\]
\end{itemize}
Then
\[
\mathbb E_\mu\!\left[
\max\bigl(\{0\}\cup\{|Z_k|:1\le k\le N\}\bigr)
\right]
\le
K_\alpha a
\sqrt{1+\max\!\left\{0,\log\!\left((b/a)^{2/\alpha}\right)\right\}}.
\]
\end{lemmav}

The two parts of the coefficient envelope have distinct statistical roles. Count feasibility caps every coefficient's stochastic scale at a constant times \(\widehat h^\beta\). The population-mass comparison and \cref{obj:lem:ordered-mass} additionally supply the rank-dependent scale
\(n^{-1/2}\widehat h^{-D_\gamma/2}k^{-1/(2(\gamma-1))}\).
The ordered maximal inequality combines these two controls. Its logarithm measures the number of ranks at which the common cap is relevant, through the scale ratio \((b/a)^{2/\alpha}\).

For the coefficient application, the decay exponent is \(1/(\gamma-1)\), and, apart from fixed multiplicative constants, the two scales are
\[
a=\widehat h^\beta,
\qquad
b=n^{-1/2}\widehat h^{-D_\gamma/2}.
\]
Their ratio is governed by the selected mesh relative to the oracle mesh. Thus the conditional maximum is controlled at the statistical scale
\(\widehat h^\beta\sqrt{1+\max\{0,\log(h_{\gamma,n}/\widehat h)\}}\).
The final inequality in \cref{obj:lem:selected-mesh} places this expression at the oracle scale even for finer selected meshes. In this application, the decay exponent ranges over
\[
\left[
\frac{1}{\gamma_{\max}-1},
\frac{1}{\gamma_{\min}-1}
\right],
\]
and the lemma supplies common constants \(K,K'\) and a common threshold \(n_0\) over the stated overlap interval. The fixed number of polynomial coordinates is included in those constants.

\subsection{The three components of the upper bound}

The approximation component uses the intrinsic Hölder regularity in \cref{obj:def:holder,obj:ass:holder-response} at degree \(m=\lceil\beta\rceil-1\). The deterministic inverse-Gram bound in \cref{obj:lem:template-conditioning} keeps the fitted conditional mean stable relative to its local polynomial approximation. On \(\mathcal E_n\), the selected-mesh comparison in \cref{obj:lem:selected-mesh} places the resulting mesh-dependent approximation error at the oracle response-recovery scale.

The conditional stochastic component is controlled by the coefficient maximum in \cref{obj:lem:selected-mesh}. On normalized partition cubes, polynomial evaluation transfers this coefficient control to spatial-supremum control, with constants determined by the fixed dimension and degree. Conditioning on the sampled design accommodates the count-based mesh selection while retaining the ordered information profile.

Finally, the clipping rule in \cref{obj:def:estimator} and the response bound in \cref{obj:def:model} bound the loss on exceptional samples by \(2B\). Consequently, the exceptional-sample contribution is at most \(2Bn^{-2}\). These three controls furnish the uniform spatial-supremum upper guarantee in \cref{obj:thm:adaptive-minimax}; the pointwise upper guarantee follows from the same spatial-supremum control.
